\addcontentsline{toc}{section}{INTRODUCCIÓN}


La adopción de la Inteligencia Artificial Generativa ha impulsado una transformación sin precedentes en la educación superior, redefiniendo las metodologías de enseñanza y la asimilación del conocimiento \parencite{Vieriu2025}. En este escenario, la gestión de apuntes universitarios ha experimentado una evolución acelerada, migrando desde los métodos manuscritos tradicionales hacia formatos digitales dinámicos e interactivos. Esta transición tecnológica busca optimizar la captura, organización y recuperación de la información académica. Según \textcite{Kreijkes2026}, la integración de Modelos de Lenguaje Grande (LLMs) presenta oportunidades excepcionales para facilitar la comprensión lectora y reducir la carga cognitiva inicial de los estudiantes. Sin embargo, el volumen exponencial de datos técnicos y científicos en disciplinas de ingeniería exige herramientas que trasciendan el simple almacenamiento, demandando ecosistemas capaces de procesar y enlazar semánticamente el contenido educativo.\\

A pesar de estos avances tecnológicos, la digitalización educativa contemporánea enfrenta desafíos algorítmicos y pedagógicos severos. Los estudiantes de ingeniería padecen de una constante sobrecarga cognitiva al intentar integrar múltiples fuentes de información no estructurada. La adopción pasiva de plataformas de IA genéricas agrava este panorama, induciendo una dependencia tecnológica que inhibe el análisis crítico y la apropiación del conocimiento \parencite{Vieriu2025}. A nivel técnico, los LLMs estándar adolecen de un problema crítico de "alucinaciones"; dado que carecen de mecanismos intrínsecos para verificar la veracidad de los hechos, frecuentemente generan información plausible pero fácticamente incorrecta \parencite{Chen2023}. Adicionalmente, la digitalización de apuntes científicos enfrenta la barrera de la notación matemática, donde los sistemas tradicionales fracasan al procesar estructuras espaciales complejas \parencite{AlvarezPerez2026}. Finalmente, las herramientas actuales promueven un consumo individualista y fragmentado de la información, careciendo de la infraestructura lógica para establecer vínculos semánticos que fomenten una colaboración auténtica entre pares \parencite{Jiang2025}.\\

Para mitigar estas brechas tecnológicas y metodológicas, la presente investigación propone formalmente el desarrollo de SynapTome, un ecosistema digital inteligente y colaborativo fundamentado en los principios de la Arquitectura Limpia (Clean Architecture). La infraestructura del sistema orquesta tecnologías de vanguardia para transformar la gestión del aprendizaje: en el núcleo de procesamiento de lenguaje natural, se implementa una arquitectura de Generación Aumentada por Recuperación (RAG) que restringe la inferencia de la IA a un corpus documental validado, garantizando precisión factual y eliminando las alucinaciones \parencite{Ning2025}. Para abordar el reto de la digitalización científica, se integra un pipeline de visión por computadora basado en el modelo Transformer TrOCR, el cual permite la transcripción estructurada de apuntes manuscritos y fórmulas matemáticas al formato LaTeX \parencite{KardanMoghaddam2025}. El elemento diferenciador de SynapTome radica en su motor de base de datos relacional modelado mediante Grafos de Conocimiento en PostgreSQL. Estos grafos indexan topológicamente los bloques de estudio, descubriendo relaciones transversales entre múltiples usuarios para catalizar una verdadera "Inteligencia Colectiva" \parencite{Obionwu2022}. Todo el ecosistema opera con una sincronización concurrente a través de SignalR y C\#, con una interfaz reactiva desarrollada en Next.js.\\

El presente trabajo de investigación se estructura en cinco capítulos secuenciales que detallan el ciclo de vida del proyecto: el Capítulo I expone el planteamiento del problema, su justificación y el contexto universitario del caso de estudio; el Capítulo II detalla los objetivos generales y específicos, así como las hipótesis de investigación y la operacionalización de variables; el Capítulo III desarrolla el marco teórico referencial, profundizando en los fundamentos de las arquitecturas RAG, TrOCR y la teoría subyacente de los Grafos de Conocimiento; el Capítulo IV describe el diseño metodológico, definiendo el enfoque experimental, las métricas de evaluación técnica y los instrumentos de usabilidad; finalmente, el Capítulo V presenta la arquitectura de software implementada, el análisis de los resultados obtenidos durante el despliegue del sistema y las conclusiones finales derivadas del proyecto.